% Pursuing Stacks (cotes 134-2 to 134-8): archive cote 134-6, site page 96 (PDF page 97), typescript page 414.
% typescript. Draft compared: ps5.tex:6421-6490.
% AI-transcribed by Claude Opus 5.5 (claude-opus-5-5, adaptive thinking, effort high), Message Batches API, batch ps3 (msgbatch_01XdFmuVVBAcgc7EFxdtCqsj), 2026-11-12; not yet checked by eye.
\scanpage{134-6/96}
\mspage{414}

As for the formula\hand{s} included in (75), the three basic ones, including
all others in a more or less formal way, are the circle\over{\unsure{s}}{d} ones a), e) and
b). T\over{\unsure{h}}{h}e properties a) and e) jointly can be viewed as the c\inserted{h}aracterizati-%
on up to canonical isomorphism, for fixed F, of the \xout{\unsure{functor}} operation
Hom(L,F), i.e. of the functor \hand{\unsure{$\widetilde{F}$}}
\[ \hand{\widetilde{F}} : \quad L \hand{\mapsto} \mathrm{Hom}(L,F) : \hand{\mathcal{A}}_{ab}{}^{o} \hand{\longrightarrow \mathcal{M}}\hand{,} \]
factoring the functor
\[ F : \hand{\mathcal{A}}^{o} \hand{\longrightarrow \mathcal{M}} \]
via the abelianization functor Wh$_{\hand{\mathcal{A}}}$ : $X \hand{\mapsto} \underline{Z}^{(X)}$ . In terms of a), the
formula b) can be viewed as essentially the definition of \underline{Hom}(L,F)\inserted[typed]{via Hom(-,F)}, more
specifically we get
\[ (76) \qquad \underline{\mathrm{Hom}}(L,F)(X) \hand{\simeq} \mathrm{Hom}(\underline{Z}^{(X)},\underline{\mathrm{Hom}}(L,F)) \hand{\simeq} \mathrm{Hom}(\underline{Z}^{(X)}\otimes L,F) \quad . \]
Taking $L = \underline{Z}^{(Y)}$ and using
\[ \underline{Z}^{(X)} \otimes \underline{Z}^{(Y)} \simeq \underline{Z}^{(XxY)} \]
((17) page 389), (76) \xout{boils down to a')} gives a'), whereas b') follows,
via (76) applied to both members,from associativity of the operation $\hand{\otimes}$ .
Formula c) is the particular case of a) for X=e, in the same way c') fol-%
lows from a'). Formula d) follows from \over{\unsure{a}}{c}),b) and the relation
\[ \underline{Z}_{\hand{\mathcal{A}}} \otimes L \simeq L \quad . \]
The exactness property e') \xout{follows} is equivalent to the similar\inserted{exactness} statement
for the functors
\[ L \hand{\mapsto} \underline{\mathrm{Hom}}(L,F)(X) \hand{\simeq} \mathrm{Hom}(\underline{Z}^{(X)}\otimes L,F) \quad , \]
for X in $\hand{\mathcal{A}}$, and thus reduces to e) with L replaced by $\underline{Z}^{(X)}\otimes L$ .

I would like now to come back to the question of ``computation'' of
Hom(L,F) and \underline{Hom}(L,F). We may for this end assume $\hand{\mathcal{A}}$ to be described by
a \uline{site} A \deleted{\illegible{1}} - which, in case the ``topology'' on A defining the site struc-%
ture is the chaotic one, brings us back to the particular case $\hand{\mathcal{A}} = A\hand{\widehat{\ }}$\hand{/}we
are mainly interested in at present. Accordingly, we'll consider the
objects F in $\hand{\mathcal{A}_{\mathcal{M}}}$ as being functors
\[ (77) \qquad \hand{F:} \ A^{o} \hand{\longrightarrow \mathcal{M}} \]
satisfying the standard exactness properties for sheaves (with respect to
the given site structure on A). In terms of (71), this is just the com-%
position of the functor (71) with the canonical functor
\[ \hand{(78)} \qquad A \hand{\longrightarrow \mathcal{A}} = A^{\hand{\sim}} , \]
associating to an object a in A the presheaf represented by it, in the

\margin{96}
